\documentclass[11pt]{article}
\usepackage[margin=1in]{geometry}
\usepackage{amsmath,amssymb,amsthm,mathtools}
\usepackage[hidelinks,pdfusetitle]{hyperref}
\usepackage{microtype,array,longtable}
\newcolumntype{L}[1]{>{\raggedright\arraybackslash}p{#1}}
\newcommand{\DOI}[1]{\mbox{\href{https://doi.org/#1}{\nolinkurl{#1}}}}
\theoremstyle{plain}
\newtheorem{theorem}{Theorem}
\newtheorem{lemma}{Lemma}
\newtheorem{proposition}{Proposition}
\newtheorem{corollary}{Corollary}
\theoremstyle{remark}
\newtheorem{remark}{Remark}
\DeclareMathOperator{\Tr}{Tr}
\DeclareMathOperator{\dist}{dist}
\newcommand{\Nodd}{N_{\mathrm{odd}}}
\newcommand{\GC}{\mathrm{GC}}
\newcommand{\src}{\mathrm{src}}
\newcommand{\phys}{\mathrm{phys}}
\setlength{\emergencystretch}{2em}
\title{Certified cubic and general-dimensional strong-coupling ceilings\newline
for electronic pair correlations in the Holstein model}
\author{Davis Rens\\SuperMatics Inc.\\\href{https://orcid.org/0000-0001-5839-5197}{ORCID 0000-0001-5839-5197}}
\date{}
\begin{document}
\maketitle

\begin{abstract}
We prove a uniform-frequency ceiling for equal-time on-site electronic pair
correlations in the spinful nearest-neighbour Holstein model on
$(\mathbb Z/N\mathbb Z)^3$. With $U=F^2/K>0$, hopping $t\ge0$, arbitrary real
common chemical potential and every phonon frequency $\Omega>0$, the bound is
$k_B T_{\mathrm{pair}}\le27900\,t^2/U$ for all $N\ge4$, improved to
$22800\,t^2/U$ for even $N$. At the corresponding thresholds, distinct-site
pair correlations are bounded by $(11/100)(99/100)^{\dist(x,y)}$.
For every integer $d\ge1$ we also prove the ceilings
$5750d^3t^2/U$ for all $N\ge4$ and $4600d^3t^2/U$ for even $N$.
The bounds transfer to translation-invariant thermodynamic canonical limits at
interior density, including balanced spin sectors. The cubic all-size constant
is about $2283$ times smaller than the earlier released hypercubic constant.
The proof combines an exact two-site trace, bridge multiplicities, proper
edge-colour H\"older decoupling, rooted-support enumeration, a degree-$2d$
Raney majorant, and a source-augmented polymer expansion. Every finite
arithmetic witness is certified by exact rational or outward-interval
arithmetic, with independently implemented checks. These are bounds on the
pair-correlation ceiling $T_{\mathrm{pair}}$, not determinations of a
superconducting transition temperature $T_c$.
\end{abstract}

\section{Model, statements and comparison}\label{sec:model}

Let $\Lambda_N^{(d)}=(\mathbb Z/N\mathbb Z)^d$, $N\ge4$, with nearest-neighbour
edges, graph distance $\dist$, and physical coordination $p=2d$.
In units $\hbar=1$, the centred spinful Holstein Hamiltonian is
\begin{equation}\label{eq:hamiltonian}
\begin{split}
H={}&-t\sum_{\langle xy\rangle,\sigma}
(c_{x\sigma}^{\dagger}c_{y\sigma}+c_{y\sigma}^{\dagger}c_{x\sigma})\\
&+\sum_x\left[\frac{P_x^2}{2M}+\frac{KQ_x^2}{2}
-FQ_x(n_x-1)-\mu_c n_x\right],
\end{split}
\end{equation}
where $n_x=n_{x\uparrow}+n_{x\downarrow}$, $K,M>0$,
$\Omega=\sqrt{K/M}>0$, $U=F^2/K>0$, and $t\ge0$.
There is no direct electron-electron interaction. The notation $\mu_c$
allows the usual harmless chemical-potential shift when an uncentred
coupling is used; throughout it is any real common chemical potential.
Write $\tau=k_BT>0$, $\beta=1/\tau$, and
$\Delta_x=c_{x\downarrow}c_{x\uparrow}$ for the physical electronic pair field.
The grand-canonical expectation is denoted by $\langle\cdot\rangle_\GC$.
The ratio $t^2/U$ has units of energy; every constant below is dimensionless.

Here $T_{\mathrm{pair}}$ is a ceiling above which distinct-site equal-time
on-site pair correlations decay exponentially and are uniformly summable.
It is not a claimed value of $T_c$. The results exclude on-site pair
 off-diagonal long-range order above their ceilings, but do not prove that
order occurs below them and do not address unrestricted extended pair fields.

\begin{theorem}[Cubic pair clustering]\label{thm:cubic}
For $d=3$, every $\Omega>0$, every real $\mu_c$, and every $N\ge4$, if
\[
\beta t^2/U\le1/27900,
\]
then, for $x\ne y$,
\begin{equation}\label{eq:cubicdecay}
|\langle\Delta_x^\dagger\Delta_y\rangle_\GC|
\le\frac{11}{100}\left(\frac{99}{100}\right)^{\dist(x,y)}.
\end{equation}
For even $N\ge4$, the same bound holds under
$\beta t^2/U\le1/22800$.
Consequently $k_BT_{\mathrm{pair}}\le27900t^2/U$ for all sizes and
$k_BT_{\mathrm{pair}}\le22800t^2/U$ for even sizes.
\end{theorem}

\begin{theorem}[A cubic-growth dimensional ceiling]\label{thm:general}
For every integer $d\ge1$ and $N\ge4$, put
\[
C_{\mathrm{all}}(d)=5750d^3,\qquad C_{\mathrm{even}}(d)=4600d^3.
\]
If $\beta t^2/U\le1/C_{\mathrm{all}}(d)$, then
\begin{equation}\label{eq:generaldecay}
|\langle\Delta_x^\dagger\Delta_y\rangle_\GC|
\le\frac{9}{40}\left(\frac{99}{100}\right)^{\dist(x,y)},\qquad x\ne y.
\end{equation}
For even $N$, $C_{\mathrm{even}}(d)$ suffices. Both assertions are uniform in
$\Omega>0$, volume and real common chemical potential, with no $t/U$ cutoff.
\end{theorem}

\begin{corollary}[Canonical thermodynamic limits]\label{cor:canonical}
The conclusions of Theorems~\ref{thm:cubic} and~\ref{thm:general} hold for
translation-invariant thermodynamic limits of total-number canonical states
at every interior density $\rho\in(0,2)$, and for balanced separately fixed
spin-number sectors at interior density. This is not a uniform finite-volume
canonical estimate, nor a claim about imbalanced separately fixed spin sectors.
\end{corollary}

For $t=0$, distinct-site pair correlations vanish; assume $t>0$ in the
estimates below. The analytical arguments are supplemented by finite
exact-rational and outward-interval computer checks. Neither those checks nor
independent scientific review amount to journal peer review or formal
proof-assistant verification.

The earlier released hypercubic bound~\cite{rensCubic},
\DOI{10.5281/zenodo.24412193}, is
\[
C_{\mathrm{old}}(d)=442368d^2(d+1)^2,
\qquad C_{\mathrm{old}}(3)=63700992.
\]
Its all-size cubic constant divided by $27900$ is exactly $1769472/775$,
approximately $2283.19$. The corresponding even-size ratio is $1327104/475$,
approximately $2793.90$. The previously released two-dimensional
result~\cite{rensSharp}, \DOI{10.5281/zenodo.24414616}, has constants $6000$
for all $N\ge4$ and $4500$ for even $N$. Those released theorems are unchanged.
The present general-dimensional bound is not claimed to improve their
specialized two-dimensional constants.

The model originates with Holstein~\cite{holstein1959}; Lang--Firsov
representations are discussed in~\cite{hohenadler2007}. Strong-coupling
expansions for attractive Holstein and Hubbard models~\cite{freericks1993}
provide qualitative context. Our argument is instead a nonperturbative
clustering estimate for the physical pair field, using the abstract polymer
criterion~\cite{kotecky1986} and correlation-expansion framework~\cite{uelschi2004}.
No numerical estimate of $T_c$, finite-frequency pair-hopping formula, or
new quantitative literature claim is used in the proof.

\section{Chronological phonon contraction and physical supports}\label{sec:contours}

The local Lang--Firsov transformation diagonalizes the atomic oscillator
coupling. After subtracting a constant, its atomic electronic Hamiltonian is
\begin{equation}\label{eq:atomic}
D_0=\frac{U}{2}\Nodd-\mu_c N_e,
\qquad \Nodd=\sum_x\boldsymbol{1}_{\{n_x=1\}}.
\end{equation}
The hopping is dressed by endpoint displacements, with displacement magnitude
squared $U/(2\Omega)$. Physical pair insertions have the corresponding
charge-two displacements. In a closed chronological Fock trajectory, the
local occupation $n_x(u)$ is periodic and piecewise constant, including
these pair jumps.

\begin{lemma}[Oscillator contraction]\label{lem:phonon}
For every such trajectory, the normalized oscillator factor is
$\exp(-\sum_x\Phi_\beta[n_x])\in(0,1]$, where
\begin{align}
\Phi_\beta[n]&=\frac14\int_0^\beta\!du\int_0^\beta\!dv\,
[n(u)-n(v)]^2K_\beta(u-v),\label{eq:phi}\\
K_\beta(w)&=\frac{U\Omega}{2}
\frac{e^{-\Omega w}+e^{-\Omega(\beta-w)}}{1-e^{-\beta\Omega}},
\qquad 0\le w\le\beta.\label{eq:kernel}
\end{align}
The kernel is periodically extended and has integral $U$.
\end{lemma}
\begin{proof}
The chronological displacement identity, including the ordering convention,
is recorded in Appendix~\ref{app:oscillator}. Integrating the kernel twice by
parts across the neutral occupation jumps gives~\eqref{eq:phi}.
The kernel is positive for every $\Omega>0$, and integration of the two
exponentials gives $\int_0^\beta K_\beta=U$. Thus $\Phi_\beta\ge0$.
Neither the number of sites nor the dimension enters this calculation.
\end{proof}

Remove the CAR signs from each legal directed spin hop and denote the
resulting symmetric, entrywise nonnegative bond matrix by $B_e$.
The chronological expansion, bounded by Lemma~\ref{lem:phonon}, is a positive
signless electronic trajectory majorant. It retains all spatial overlaps and
all hopping orders; no factorization in time or restriction to initially
even occupations is made.

Group terms by their distinct used edge labels $E$. Components with disjoint
vertex sets factor, since their atomic traces are independent and their
fermion-even bond or pair-source operators commute on disjoint sets.
Chronological shuffles then give a hard-core spatial polymer gas with
connected, nonempty edge supports and incompatibility by vertex intersection.
Write $w(E)$ for a physical activity and $w_{xy}(E)$ for a pair-marked activity.
The usual connected-support identity is
\begin{equation}\label{eq:pairpolymer}
\langle\Delta_x^\dagger\Delta_y\rangle_\GC
=\sum_{\substack{E\text{ connected}\\x,y\in V(E)}}
w_{xy}(E)\frac{Z_{\Lambda\setminus V(E)}}{Z_\Lambda}.
\end{equation}
The dimension-independent support representation, bridge and trace lemmas
are the same local lemmas used in the released two-dimensional
proof~\cite{rensSharp}, Sections 3--6 and its oscillator appendix.
We repeat their quantitative content and derive all geometric dependence below.

\section{The exact bond and its fugacity supremum}\label{sec:bond}

Put $b=\beta U$, $q_b=e^{-b/2}$, $z=e^{\beta\mu_c}$, and let $g\ge0$ be the
dimensionless hopping argument. The one-site atomic partition polynomial is
$z_a(z)=1+2q_bz+z^2$.

\begin{lemma}[Two-site spectrum and uniform fugacity bound]\label{lem:bond}
The normalized two-site trace ratio is
\begin{equation}\label{eq:exactratio}
R(b,g,z)=\frac{1+z^4+4q_b(z+z^3)\cosh g
+z^2(1+3q_b^2+2q_b\cosh r)}{z_a(z)^2},
\quad r=\sqrt{b^2/4+4g^2}.
\end{equation}
It is maximized over $z>0$ at $z=1$, and
\begin{equation}\label{eq:bondub}
R(b,g,z)\le R_{\mathrm{ub}}(b,g):=
\frac{3+4q_b^2+8q_b\cosh g+e^{r-b/2}}{4(1+q_b)^2}.
\end{equation}
\end{lemma}
\begin{proof}
The two-site Fock space has dimension 16. In the zero- and four-electron
sectors the electronic exponent is zero. Each of the one- and three-electron
sectors has eigenvalues $-b/2\pm g$, each with multiplicity two.
In the two-electron sector they are $-b$ with multiplicity three, $0$ with
multiplicity one, and $-b/2\pm r$. Tracing these five sectors gives
\eqref{eq:exactratio}. This is also a direct diagonalization of the
signless single-bond matrix, which is gauge-equivalent to its two-site CAR matrix.

Let $w=z+z^{-1}\ge2$. Subtraction of the atomic denominator gives
\[
R=1+\frac{Aw+B_0}{(w+2q_b)^2},\quad
A=4q_b(\cosh g-1),\quad B_0=2q_b\cosh r-q_b^2-1.
\]
Here $A\ge0$ and $B_0\ge0$, because $r\ge b/2$.
The derivative in $w$ has numerator $-Aw+2Aq_b-2B_0\le0$.
Thus $w=2$, equivalently $z=1$, is worst, uniformly in $\mu_c$.
Finally $e^{-r-b/2}\le e^{-b}=q_b^2$ gives~\eqref{eq:bondub}.
The coefficient $3q_b^2$ belongs to the exact spectrum; $4q_b^2$ appears
only after this upper bound. No atomic fugacity penalty is introduced.
\end{proof}

\section{Bridge multiplicities, colours and the free scale}\label{sec:weights}

For connected $E$, write $k=|E|$, $v=|V(E)|$, $h(E)$ for its bridge count,
and $m=k+h(E)$. A physical edge is active only if at least one hop uses it.

\begin{lemma}[Multiplicity and coloured trace bound]\label{lem:weight}
Every active physical bridge has at least two hops, both in a vacuum contour
and in a pair-marked contour. Thus the total physical hopping multiplicity is
at least $m$. The same assertion holds within each connected physical
component of a source-augmented contour.
If $E$ has a proper $c$-edge-colouring and $Cy\ge1$, define
\[
\ell=\sqrt{Cy},\qquad
M(c,y)\ge\sup_{b\ge0}R_{\mathrm{ub}}(b,c\sqrt{yb}).
\]
Under $\beta t^2/U\le1/C$,
\begin{equation}\label{eq:weight}
|w(E)|,\ |w_{xy}(E)|\le
R_E:=(Cy)^{-m/2}M(c,y)^{k/c}.
\end{equation}
\end{lemma}
\begin{proof}
Removing a bridge partitions the physical component into two sides.
A hop through it changes the electron number on one side by one.
Closed trajectories return that number modulo two; pair insertions and pair
transfers change it by zero or two. The bridge hopping count is therefore
even and, for an active edge, at least two. Nonbridges require at least one
hop, proving the lower bound $m$.

Scale all physical hopping labels by $\ell\ge1$. Each trajectory's scale
penalty is at most $\ell^{-m}$. Apply Schatten H\"older, or log-trace
convexity together with the product formula~\cite{trotter1959}, to the $c$
matching colour classes. It multiplies the hopping argument on each matching
by $c$. The matching trace factors into single-bond traces and untouched
atomic sites. The pair insertion has operator norm at most one and its
endpoint trace is bounded by the same positive contour argument.
These dimension-independent trace statements are the coloured H\"older and
pair norm lemmas of~\cite{rensSharp}, Section 6. The product of the single-bond
ratios, to the power $1/c$, gives $M^{k/c}$.

Indeed, with $b=\beta U$ and $g=c\ell\beta t$,
\[
g^2=c^2Cy\,\beta U\,(\beta t^2/U)\le c^2yb.
\]
The bond ratio increases with $g\ge0$, so the stated supremum suffices.
This proves~\eqref{eq:weight} without a restriction on $\Omega$ or $\mu_c$.
\end{proof}

The scale is $y=\ell^2/C$, not $\ell^2$. Its admissibility is the lower
constraint $Cy\ge1$. At fixed $y>0$ the finite bond constant is independent
of $C$, and the tail edge weight
\begin{equation}\label{eq:tailweight}
w=\frac{B M(c,y)^{1/c}}{q\sqrt{Cy}}
\end{equation}
tends to zero as $C$ increases. The factors $B>1$ and $q\in(0,1)$ will be
the vertex budget and decay penalty. A class-specific choice of $y$ is
allowed: it bounds each support class separately and does not alter the
Hamiltonian. Choosing an unproved smaller colour count is not allowed.

\begin{lemma}[Proper colour counts]\label{lem:colours}
The torus has edge chromatic number $2d$ for even $N$ and $2d+1$ for odd $N$.
For a small simple support $E$, a valid cap is its maximum degree $\delta(E)$
if it is bipartite, and $\delta(E)+1$ otherwise.
\end{lemma}
\begin{proof}
For even $N$, direction $j$ uses two colours distinguished by the parity
of the starting coordinate $x_j$. They form $2d$ matchings; degree gives
the lower bound.
For an odd cycle alternate two colours on its non-wrap edges and put a
third colour on its wrap edge. Every vertex has exactly one missing colour.
Inductively copy a $(2d-1)$-colouring of the $(d-1)$-torus into each layer,
use two new colours alternately on the new non-wrap edges, and colour the
new wrap edge through old vertex $a$ by the unique old colour missing at $a$.
That colour is missing at both endpoints, and the wrap edges form a matching.
This gives $2d+1$ colours. For odd $N$, a matching has at most
$(N^d-1)/2$ edges; $2d$ such matchings cannot cover the $dN^d$ physical edges.
The class caps are the bipartite edge-colouring theorem and the simple-graph
maximum-degree-plus-one theorem. A three-edge claw in particular needs three
colours. The physical degree remains $p=2d$ even on an odd torus.
\end{proof}

\section{Rooted support counts and the cubic class table}\label{sec:counting}

Fix an ordering of incident physical edges at each vertex. Deterministic
labelled depth-first exploration encodes a connected edge set in $p$ root
slots and $p-1$ forward slots per physical edge token. An edge to an already
visited vertex is a terminal token; its ambient slot identifies its endpoint.
The code reconstructs the support, including cycles, so the exploration is
injective into the ordered-tree majorant.

\begin{lemma}[Degree-$p$ Raney majorant]\label{lem:raney}
Let the least nonnegative solution satisfy
\begin{equation}\label{eq:raney}
T=1+wT^{p-1},\qquad G_{\phys}=T^p.
\end{equation}
The number of connected rooted $k$-edge supports is at most
\begin{equation}\label{eq:raneycoeff}
R_p(k)=[w^k]T^p=
\frac{p}{(p-1)k+p}\binom{(p-1)k+p}{k}.
\end{equation}
For rational $T_0>1$ with $(p-1)(T_0-1)/T_0<1$, put
$w_{\mathrm{eff}}=(T_0-1)/T_0^{p-1}$.
For $0\le w\le w_{\mathrm{eff}}$, the $k\ge7$ tail is bounded by
\begin{equation}\label{eq:raneytail}
T_0^p-\sum_{k=0}^6 R_p(k)w_{\mathrm{eff}}^k.
\end{equation}
\end{lemma}
\begin{proof}
The preceding slot encoding gives~\eqref{eq:raney}; Lagrange inversion gives
\eqref{eq:raneycoeff}. For $p>2$ the branch condition selects the small root
before its critical point $(p-1)/(p-2)$. For $p=2$, $T=(1-w)^{-1}$ and
$w_{\mathrm{eff}}<1$. The tail series has nonnegative coefficients and is
increasing in $w$, so its complete value at $w_{\mathrm{eff}}$ is an upper
bound. The root and the subtracted polynomial are evaluated at the same
weight; no mixed-weight tail subtraction is used.
\end{proof}

For $p=6$, the coefficients through $k=6$ are
\[
1,\ 6,\ 45,\ 380,\ 3450,\ 32886,\ 324632.
\]
They are independently re-derived by formal polynomial fixed-point iteration
in the certificate, without using the binomial formula.

\begin{proposition}[Cubic small-support domination]\label{prop:classes}
Table~\ref{tab:classes} dominates rooted connected supports with $k\le6$ on
all even cubic tori $N\ge4$. To dominate every $N\ge4$, add the two starred
rows. The row fields are $(v,k,c,m,n)$, where $n$ is the rooted count and $c$
is the valid cap of Lemma~\ref{lem:colours}. The even table has 21 rows and
count sum $337498$; the all-size table has 23 rows and sum $337573$.
\end{proposition}

\begin{table}[htbp]
\centering\small
\begin{tabular}{rrrrr@{\qquad}l}
$v$&$k$&$c$&$m$&$n$&\text{scope}\\\hline
2&1&1&2&6&both\\
3&2&2&4&45&both\\
4&3&2&6&300&both\\
4&3&3&6&80&both\\
4&4&2&4&15&both\\
5&4&2&8&1815&both\\
5&4&3&8&1500&both\\
5&4&4&8&75&both\\
5&5&3&5&3&all sizes only ($*$)\\
5&5&3&6&300&both\\
6&5&2&10&10602&both\\
6&5&3&10&18936&both\\
6&5&4&10&1800&both\\
6&5&5&10&36&both\\
6&6&2&6&240&both\\
6&6&3&8&3960&both\\
6&6&4&7&72&all sizes only ($*$)\\
6&6&4&8&540&both\\
7&6&2&12&59241&both\\
7&6&3&12&208012&both\\
7&6&4&12&28938&both\\
7&6&5&12&1050&both\\
7&6&6&12&7&both\\\hline
\end{tabular}
\caption{Classwise-max cubic support counts.}\label{tab:classes}
\end{table}

\begin{proof}[Proof of Proposition~\ref{prop:classes}]
Two complete enumerators agree class-by-class at $N=4,5,6,7$, with respective
totals $329734$, $335959$, $336409$, $336427$ and class counts $21$, $23$, $21$,
$21$. The primary uses levelwise integer edge masks, Tarjan low-link bridges
and breadth-first bipartiteness. The secondary uses reversed edge ordering,
frozenset supports, edge-deletion reachability for bridges and vertex-incidence
row reduction over $\mathbb F_2$ for bipartiteness. Each retains each support
once. Taking rowwise maxima over $N=4,6,7$ gives the even table; adding $N=5$
gives the all-size table. The count sum of a rowwise maximum is not claimed
to be the count on a single torus.

For $N\ge7$, every cycle in a $k\le6$ support has length below $N$, so cannot
wind around any coordinate. Lifting along a spanning tree gives signed unit
coordinate displacements with zero displacement around all cycles. Distinct
vertices cannot collide because their connecting tree path has at most six
edges, too few for a nonzero displacement multiple of $N$. Conversely,
reduction modulo $N$ embeds each such rooted lattice support. Thus the table
stabilizes for all $N\ge7$. The certificate additionally checks $T_7=T_8$.
It also reproduces the independent two-dimensional 18/20 class guards from
the released proof~\cite{rensSharp}, Section 7.
\end{proof}

\section{Certified bond pins and the cubic vertex sum}\label{sec:cubicbudget}

Put $a=c^2y$ and
\[
h(b)=\sqrt{b^2/4+4ab}-b/2
=\frac{4ab}{\sqrt{b^2/4+4ab}+b/2},\qquad h(0)=0.
\]
This function increases from zero to $4a$. On a closed interval
$[b_0,b_1]$, bound $q_b$ between its endpoint values, $g$ above by
$\sqrt{ab_1}$, and $h$ above by its value at $b_1$, using a lower square-root
enclosure in the rationalized expression. Bounding the numerator of
\eqref{eq:bondub} upwards and its denominator downwards gives a valid
interval bound, not a presumption that a grid locates the supremum.

For $b\ge40$ all the chosen pins have $40\ge4a$. Hence
$-b/2+\sqrt{ab}$ is decreasing, $h\le4a$, and
\begin{equation}\label{eq:inftail}
R_{\mathrm{ub}}\le
\frac{3+4e^{-40}+4e^{-20+\sqrt{40a}}+4e^{-20}+e^{4a}}4.
\end{equation}
Adaptive closed-interval bisection covers $[0,40]$ and this formula covers
$[40,\infty)$. Both the $b=0$ endpoint and the infinite tail are included.

\begin{proposition}[Uniform bond pins]\label{prop:pins}
For every row of Table~\ref{tab:pins}, its rational $M$ bounds
$\sup_{b\ge0}R_{\mathrm{ub}}(b,c\sqrt{yb})$.
\end{proposition}
\begin{table}[htbp]
\centering\small
\begin{tabular}{rrrr}
$c$&$m/k$&$y$&$M$\\\hline
1&2&$299/1000$&$79/50$\\
2&2&$39/250$&$1893/500$\\
2&1&$1/10$&$249/125$\\
3&2&$23/250$&$1903/250$\\
3&1&$57/1000$&$1349/500$\\
3&$6/5$&$13/200$&$837/250$\\
3&$4/3$&$69/1000$&$15/4$\\
4&2&$13/200$&$1677/100$\\
4&$7/6$&$43/1000$&$584/125$\\
4&$4/3$&$6/125$&$123/20$\\
5&2&$51/1000$&$20879/500$\\
6&2&$21/500$&$10657/100$\\
6&1 (tail)&$23/1000$&$1903/250$\\
7&1 (tail)&$19/1000$&$1111/100$\\\hline
\end{tabular}
\caption{Bond bounds; $m/k$ selects the shared class pin.}\label{tab:pins}
\end{table}
\begin{proof}
The first implementation uses outward 48-bit dyadic rational rounding,
Taylor/range-reduced exponentials, rounded Newton upper square roots and
reciprocal lower square-root bounds. The independent implementation uses an
unrounded 24-term rational Taylor sum with a positive geometric remainder,
rational powers, and integer-square-root enclosures. They each cover the
14 distinct pins on 101 accepted cells, plus~\eqref{eq:inftail}.
All branch and enclosure decisions are exact integer or rational comparisons.
Section~\ref{sec:certificate} describes the arithmetic independently of its
floating-point display output.
\end{proof}

For $q=99/100$, define a rooted physical vertex upper bound
\begin{equation}\label{eq:physicalbudget}
\begin{split}
S_{\phys}={}&\sum_{k\le6\text{ classes}}
 n B^v (Cy)^{-m/2}M^{k/c}q^{-k}\\
&+B\left[T_0^6-\sum_{k=0}^6R_6(k)w_{\mathrm{eff}}^k\right].
\end{split}
\end{equation}
For the tail, $v\le k+1$ and $m\ge k$ give the outside factor $B$ and the
per-edge factor~\eqref{eq:tailweight}. In particular, $B$ is not omitted
from that edge factor. We impose the stronger source-compatible domination
in the next section, which implies $w\le w_{\mathrm{eff}}$.

Fractional powers in the finite sum are bounded by rational upper values
$R=j/10^{10}$ satisfying
\begin{equation}\label{eq:rootcomparison}
R^{2c}(Cy)^{mc}\ge M^{2k}.
\end{equation}
This is an integer-power comparison, not numerical root extraction.
The rowwise scales all satisfy $Cy\ge1$.

\section{Source multigraph and the polymer criterion}\label{sec:source}

Fix a nonzero displacement $r$ and phase $\theta$, and add the Hermitian,
number-conserving source
\[
X_r=\sum_x[e^{i\theta}\Delta_x^\dagger\Delta_{x+r}+\mathrm{h.c.}],
\qquad H(\xi)=H-\xi X_r.
\]
It has at most two labelled source incidences per vertex, even when
$r=-r$ modulo $N$, in which case parallel labels form a multigraph.
There is no dependence of this source degree on $d$ or $|r|$.
Set
\begin{equation}\label{eq:sourcedisk}
L=100000,\qquad \xi_0=\frac{\tau}{2L}=\frac{\tau}{200000},
\qquad |\xi|\le\xi_0.
\end{equation}
The chronological norm estimate for $k_s$ source labels and connected
physical components $E_i$ is
\begin{equation}\label{eq:sourceactivity}
|w(E_\phys\sqcup E_\src)|
\le L^{-k_s}(e^{2\beta L|\xi|}-1)^{k_s}\prod_i R_{E_i}.
\end{equation}
To obtain it, separate source insertion times, bound the normalized phonon
factor by one, and use Schatten H\"older on the intervening positive
propagators $e^{\Delta s H_\phys}$, whose times sum to $\beta$, with the
source operator norm at most one. This is the dimension-independent
source estimate of the released two-dimensional proof~\cite{rensSharp},
Section 10. On~\eqref{eq:sourcedisk}, $e^{2\beta L|\xi|}-1\le e-1<2<3$;
therefore $w_s=3B/L$ bounds a source edge after the vertex factor is absorbed.
Contracting physical components gives $v\le k_p+k_s+1$.

\begin{lemma}[Source-compatible exploration]\label{lem:source}
The labelled physical/source exploration has generating functions
\begin{equation}\label{eq:twovar}
P=1+w_pP^{p-1}S^2,\qquad S=(1-w_sP^p)^{-1},\qquad G=P^pS^2.
\end{equation}
For $T=T_0$, define
\begin{equation}\label{eq:sourceparams}
S_\mathrm{hi}=(1-w_sT^p)^{-1},\qquad
w_p=\frac{T-1}{T^{p-1}S_\mathrm{hi}^2}.
\end{equation}
If $w_sT^p<1$ and
\begin{equation}\label{eq:den}
\mathrm{den}=1-(p-1)w_pT^{p-2}S_\mathrm{hi}^2
-2p w_pw_sT^{2p-2}S_\mathrm{hi}^3>0,
\end{equation}
then a rooted source-containing vertex upper bound is
\begin{align}
D_P&=\frac{2w_pT^{2p-1}S_\mathrm{hi}^3}{\mathrm{den}},\label{eq:Dp}\\
D_S&=T^pS_\mathrm{hi}^2+pw_sT^{p-1}S_\mathrm{hi}^2D_P,\label{eq:Ds}\\
S_\src&=Bw_s[pT^{p-1}S_\mathrm{hi}^2D_P+2T^pS_\mathrm{hi}D_S].\label{eq:srcbudget}
\end{align}
It applies whenever, for every physical component,
\begin{equation}\label{eq:component}
B^kR_Eq^{-k}\le w_p^k.
\end{equation}
\end{lemma}
\begin{proof}
A physical arrival leaves $p-1$ physical slots and two source slots; a source
arrival leaves $p$ physical slots and one forward source slot.
This gives~\eqref{eq:twovar}, also when labels are parallel. The rectangle
$1\le P\le T$, $1\le S\le S_\mathrm{hi}$ is invariant under monotone iteration
from $(1,1)$. Differentiating both implicit equations, including the dependence
of $S$ on $P$, gives~\eqref{eq:den}--\eqref{eq:Ds}.
The term $2p=4d$, equal to 12 at $d=3$, must be retained.
All series coefficients are nonnegative, so evaluation at the rectangle
boundary bounds their derivatives. The sum over at least one source label is
at most $w_s\partial_{w_s}G$, yielding~\eqref{eq:srcbudget}.
Equation~\eqref{eq:component} and~\eqref{eq:sourceactivity} then justify
assigning each physical token weight $w_p$.
\end{proof}

For the $k\ge7$ components, a sufficient exact rational form of
\eqref{eq:component} is
\begin{equation}\label{eq:tailsource}
(B^cM_{\mathrm{tail}})^2
\le(w_pq)^{2c}(Cy_{\mathrm{tail}})^c.
\end{equation}
For every small row it is checked with~\eqref{eq:rootcomparison}.
The independent eliminated Jacobian identity is
\[
\mathrm{den}=1-(p-1)(T-1)/T
-2p(T-1)w_sT^{p-1}S_\mathrm{hi}.
\]

\begin{lemma}[KP deletion and pair bound]\label{lem:kp}
If $S_\phys+S_\src\le\ln B$, then the source expansion is analytic on
\eqref{eq:sourcedisk}, the deletion ratio satisfies
$|Z_{\Lambda\setminus S}/Z_\Lambda|\le B^{|S|}$, and at zero source
\[
|\langle\Delta_x^\dagger\Delta_y\rangle_\GC|
\le(\ln B)q^{\dist(x,y)}.
\]
\end{lemma}
\begin{proof}
Use the Koteck\'y--Preiss criterion~\cite{kotecky1986} with
$a(E)=|V(E)|\ln B$ and the additional physical-edge decay penalty
$b(E)=|E_\phys|\ln(1/q)$. A vertex sum bounded by $\ln B$ implies the
incompatibility sum by a union bound over the vertices of the test polymer.
Pinning each vertex to be deleted gives the stated deletion bound.
In~\eqref{eq:pairpolymer}, each support connecting $x,y$ has at least
$\dist(x,y)$ physical edges. Extracting its decay factor gives the result.
The source estimate is used for analyticity, not for a claimed nonzero-source
pair decay inequality.
\end{proof}

\begin{proposition}[Cubic joint KP witnesses]\label{prop:cubicbudget}
For $q=99/100$ and $L=100000$, the following rational choices satisfy all
small-row and tail scale, source-component, small-branch, pole, Jacobian,
vertex-budget and prefactor inequalities:
\[
\begin{array}{c|ccc}
\text{scope}&C&B&T_0\\\hline
\text{even }N&22800&277/250&449/400\\
\text{all }N&27900&111/100&1123/1000
\end{array}
\]
The conservative decimal summaries are
\[
\begin{array}{c|ccc}
& S_\phys\text{ upper}&S_\src\text{ upper}&\ln B\text{ lower}\\\hline
\text{even}&0.101751108&0.000719841&0.102556588\\
\text{all}&0.103493088&0.000729760&0.104360015
\end{array}
\]
with positive margins exceeding $0.000085641$ and $0.000137168$,
respectively. Both upper bounds on $\ln B$ are below $11/100$.
\end{proposition}
\begin{proof}
The certificate computes the finite sums, tails and source bounds using
exact rational comparisons and the enclosures in Proposition~\ref{prop:pins}.
The displayed decimals are conservatively rounded summaries, not proof inputs.
Its independent verification re-certifies all bond pins and all small-support
counts. Applying Lemma~\ref{lem:kp} proves Theorem~\ref{thm:cubic}.
\end{proof}

\section{A uniform cubic-growth bound in every dimension}\label{sec:general}

The general-dimensional argument separates the single-edge support and uses
all other supports through the $k\ge2$ Raney tail, so no small-support table
is needed.

\begin{lemma}[A simple uniform bond estimate]\label{lem:simplebond}
For $a=c^2y$ and all $b\ge0$,
\begin{equation}\label{eq:simplebond}
R_{\mathrm{ub}}(b,c\sqrt{yb})\le\frac34+e^{a/2}+\frac14e^{4a}.
\end{equation}
For $c\ge4$ and $y=1/(8c)$, the right side is at most $e^{c/2}$.
Thus, whenever $Cy\ge1$ and $m\ge k$,
\begin{equation}\label{eq:generalweight}
R_E\le\left(\sqrt{8ec/C}\right)^k.
\end{equation}
For the one-edge support, use $c=1$, $m=2$, $y=1/8$ and
\begin{equation}\label{eq:onebond}
M_1=\frac34+e^{1/16}+\frac14e^{1/2}<\frac{23}{10},
\qquad R_1\le8M_1/C.
\end{equation}
\end{lemma}
\begin{proof}
In~\eqref{eq:bondub}, $(3+4q_b^2)/(4(1+q_b)^2)\le3/4$ and $h\le4a$.
Young's inequality gives $\sqrt{ab}\le(a+b)/2$, hence
$\cosh\sqrt{ab}\le e^{a/2}\cosh(b/2)$.
Also $2q_b\cosh(b/2)/(1+q_b)^2=(1+q_b^2)/(1+q_b)^2\le1$.
This proves~\eqref{eq:simplebond}.
For $a=c/8\ge1/2$, put $u=e^{a/2}>5/4$.
The polynomial $u^8-(4/3)u-1$ is positive at $5/4$ and increasing thereafter,
so $3/4+e^{a/2}+e^{4a}/4\le e^{4a}=e^{c/2}$.
The two rational polynomial checks and the Taylor lower bound
$e^{1/4}>5/4$ are explicit certificate checks.
Lemma~\ref{lem:weight} now gives~\eqref{eq:generalweight}.
Applying~\eqref{eq:simplebond} to the one-edge parameters gives
\eqref{eq:onebond}; its rational Taylor upper bound is also certified.
The special $c\ge4$ estimate is not claimed for smaller $c$.
\end{proof}

\subsection{Dimension-free rational reduction for \texorpdfstring{$d\ge2$}{d >= 2}}\label{sec:generalhigh}

Choose
\begin{equation}\label{eq:generalparams}
B=5/4,\quad \alpha=7/20,\quad q=99/100,\quad L=100000,
\quad p=2d,\quad T=1+\alpha/p.
\end{equation}
Let rational Taylor upper bounds satisfy $E_\alpha\ge e^\alpha$ and $E_1\ge e$,
and put
\[
w_s=3B/L,\quad \overline S=(1-w_sE_\alpha)^{-1},\quad
\overline D=1-\alpha-2\alpha w_sE_\alpha\overline S.
\]
The certificate verifies $w_sE_\alpha<1$ and $\overline D>0$.
Since $T^p\le e^\alpha\le E_\alpha$ and
$(p-1)(T-1)/T<\alpha<1$, this is on the small physical branch,
$S_\mathrm{hi}\le\overline S$, $\mathrm{den}\ge\overline D$, and
\begin{equation}\label{eq:wpbound}
w_p\ge\frac{\alpha}{pE_\alpha\overline S^2}.
\end{equation}

Put
\begin{equation}\label{eq:base}
\mathcal B=E_1\left(\frac{B\overline S^2E_\alpha}{q\alpha}\right)^2.
\end{equation}
The source-compatible tail condition follows from~\eqref{eq:generalweight}
and~\eqref{eq:wpbound} if $C\ge32cd^2\mathcal B$.
For even sizes $c=2d$, this is $C\ge64d^3\mathcal B$.
For all sizes, $c\le2d+1$ and $(2d+1)/d\le5/2$ when $d\ge2$, so
$C\ge80d^3\mathcal B$ suffices.
The exact rational upper evaluations satisfy
\begin{equation}\label{eq:Kvalues}
64\mathcal B<4560.229<4600,\qquad
80\mathcal B<5700.286<5750.
\end{equation}

Scale admissibility is explicit: for $d\ge2$ the even tail has
$Cy=4600d^2/16\ge1150$, and the all-size tail has
$Cy=5750d^3/[8(2d+1)]\ge1150$, because this ratio increases with $d$.
The one-edge scale has $Cy=C/8\ge4600/8=575$ for every $d\ge1$.
The certificate checks the two $d=2$ tail minima and this global one-edge
minimum. Thus the use of $m\ge k$ in~\eqref{eq:generalweight} is admissible.

The $k\ge2$ rooted physical sum is bounded by
\begin{equation}\label{eq:generaltail}
S_{\mathrm{tail}}\le B[E_\alpha-1-\alpha/E_\alpha],
\end{equation}
because $pw_{\mathrm{eff}}=\alpha/T^{p-1}\ge\alpha/E_\alpha$.
The number of one-edge rooted supports is $p=2d$, not $d$, giving
\begin{equation}\label{eq:generalone}
S_{\mathrm{one}}\le\frac{16B^2M_1}{4600d^2q}
\le\frac{4B^2M_1}{4600q},\qquad d\ge2.
\end{equation}
For its source-component domination, a sufficient worst-case $d=2$ condition is
\begin{equation}\label{eq:onesource}
\frac{4BM_1}{4600q}\le\frac{\alpha}{E_\alpha\overline S^2}.
\end{equation}

Define
\begin{align}
J&=2\alpha E_\alpha^2\overline S^3/\overline D,\label{eq:J}\\
\overline D_S&=E_\alpha\overline S^2+w_sJ,\label{eq:DbarS}\\
S_{\src}&\le Bw_s[J+2E_\alpha\overline S\,\overline D_S].\label{eq:generalsource}
\end{align}
These follow directly from~\eqref{eq:Dp}--\eqref{eq:srcbudget}: substitute
$w_p=(T-1)/(T^{p-1}S_\mathrm{hi}^2)$, use $p(T-1)=\alpha$,
$T^{2p-1}\le E_\alpha^2$, and $\mathrm{den}\ge\overline D$.
In particular $pT^{p-1}S_\mathrm{hi}^2D_P\le J$, which also bounds the
implicit contribution in $D_S$. The source Jacobian factor $2p$ is retained.

\begin{proposition}[Dimension-free joint budget]\label{prop:generalbudget}
The rational upper bounds in~\eqref{eq:generalparams}--\eqref{eq:generalsource}
satisfy~\eqref{eq:Kvalues},~\eqref{eq:onesource}, all scale constraints, and
\[
\begin{aligned}
S_{\mathrm{tail}}&<0.215533397,&
S_{\mathrm{one}}&<0.003055933,\\
S_{\src}&<0.000290509,&
\ln(5/4)&>0.223143551.
\end{aligned}
\]
The joint KP margin exceeds $0.004263714$ and $\ln(5/4)<9/40$.
\end{proposition}
\begin{proof}
Unrounded rational Taylor/geometric upper bounds on the exponentials and
rational lower/upper logarithm bounds certify these inequalities.
Independent finite-$d$ evaluations of the exact source equations at
$d=2,3,4,7,100$ additionally check the directions of the eliminated source
and tail bounds. The all-$d$ proof is the displayed majorant, not extrapolation
from those sample dimensions. Lemma~\ref{lem:kp} proves
Theorem~\ref{thm:general} for $d\ge2$.
\end{proof}

\subsection{Dimension one, without the \texorpdfstring{$c\ge4$}{c >= 4} shortcut}\label{sec:d1}

For $d=1$, use $p=2$, $T=1+\alpha/2=47/40$ and the exact source equations
\eqref{eq:sourceparams}--\eqref{eq:srcbudget}. Let $(c,K)$ be $(2,4600)$ for
even sizes and $(3,5750)$ for all sizes. Set
\[
y=1/(8c),\qquad M_c=3/4+e^{c/16}+e^{c/2}/4.
\]
Use~\eqref{eq:simplebond}, not the $c\ge4$ shortcut. The exact required
inequalities are
\begin{align}
(B^cM_c)^2&\le(w_pq)^{2c}(Ky)^c,\label{eq:d1tail}\\
8BM_1/(Kq)&\le w_p,\label{eq:d1one}\\
B[T^2-1-2(T-1)/T]+16B^2M_1/(Kq)+S_{\src}&\le\ln B.\label{eq:d1kp}
\end{align}
Scale admissibility is $Ky=575/2$ in the even case and $2875/12$ in the
all-size case, both above one; $K/8\ge575$ for the one-edge scale.
All inequalities, including $Ky\ge1$, are exact certificate guards.
The joint sums are below $0.115906$ and $0.113462$, respectively, while
$\ln B>0.223143551$.
This completes Theorem~\ref{thm:general} for every $d\ge1$.

\section{Canonical transfer and uniform summability}\label{sec:canonical}

We give the pressure argument to specify the ensemble scope and the exact
Cauchy radius. It is the canonical route used in the released
proof~\cite{rensSharp}, Section 10, with its dimension dependence made explicit.
For each fixed displacement $r$, the source expansion is analytic on
$|\xi|\le\xi_0=\tau/200000$. The absolute cluster budget per site is below
one, so its normalized log partition difference from zero source is bounded
by two. Cauchy's estimate on the exact half-radius disk
\begin{equation}\label{eq:cauchydisk}
|\xi|\le\xi_0/2=\tau/400000
\end{equation}
bounds the quadratic remainder by $4\xi^2/\xi_0^2$ for real $\xi$.
The derivative at zero source is
$2\beta\operatorname{Re}(e^{i\theta}\langle\Delta_0^\dagger\Delta_r\rangle_\GC)$.
Thus, with the appropriate $A$ and $q$ from the main theorem,
\begin{equation}\label{eq:pressure}
\frac1V\log\frac{Z_\GC(\xi;\mu_c)}{Z_\GC(0;\mu_c)}
\le2\beta A q^{\dist(0,r)}|\xi|+4\xi^2/\xi_0^2.
\end{equation}
This uses zero-source pair decay only. Analyticity at nonzero source is not
mistaken for a source-augmented pair decay theorem.

Divide a large $d$-torus into fixed open blocks of side $L_b$. For each fixed
$d$ and $r$, the dropped hopping and source boundary terms are bounded and
have boundary-to-volume cost $O(1/L_b)$. The disconnected canonical partition
is a sum over block number allocations; its log differs from the largest
allocation log by at most the number of blocks times $\ln(2L_b^d+1)$.
The per-site allocation entropy is therefore at most
\[
L_b^{-d}\ln(2L_b^d+1)\longrightarrow0.
\]
Atomic oscillator sector traces are finite, and electron occupations per site
are finite. The block-subdivision equivalence-of-ensembles argument extends
verbatim to $(\mathbb Z/N\mathbb Z)^d$ for each fixed $d$: the boundary and
allocation terms vanish, and spin exchange is dimension-independent.
It gives a finite concave canonical density function and its grand-canonical
pressure dual. At an interior total density choose a finite supporting common
chemical potential. Along the canonical thermodynamic sequence,
\[
\delta_V=V^{-1}[\log Z_\GC(0;\mu_c)-\log Z_{N_e}(0)-\beta\mu_cN_e]
\ge0,\qquad \delta_V\longrightarrow0.
\]
Since $e^{\beta\mu_cN_e}Z_{N_e}(\xi)\le Z_\GC(\xi;\mu_c)$ for real $\xi$,
combine this fact with~\eqref{eq:pressure}. Convexity of the canonical
real-source log partition bounds its derivative by secants, giving
\begin{equation}\label{eq:canonicalsecant}
\begin{split}
|2\beta\operatorname{Re}(e^{i\theta}
\langle\Delta_0^\dagger\Delta_r\rangle_{N_e})|
\le{}&2\beta A q^{\dist(0,r)}\\
&+4\xi/\xi_0^2+\delta_V/\xi,
\quad 0<\xi\le\xi_0/2.
\end{split}
\end{equation}
Take the thermodynamic limit first at fixed $\xi$, then $\xi\to0$, and vary
$\theta$. This proves the stated canonical bound for each fixed displacement.
It does not provide a uniform finite-volume canonical deletion ratio.

For separately fixed spin sectors, use allocations of $(N_\uparrow,N_\downarrow)$.
Their entropy per block is at most $2\ln(L_b^d+1)$, still vanishing per site.
At a balanced interior density, a finite supporting chemical-potential vector
and its spin-swapped vector are both supporting, so their average has equal
components. A real common chemical potential thus supports the balanced
point. The source preserves both spin numbers, and the same secant proof
applies. No common-potential support is asserted at imbalanced densities.
This proves Corollary~\ref{cor:canonical}.

For fixed $d$ the number of lattice sites at distance $s\ge1$ is
\begin{equation}\label{eq:shell}
N_d(s)=\sum_{j=1}^{\min(d,s)}2^j\binom dj\binom{s-1}{j-1}.
\end{equation}
Choose the nonzero coordinates, their signs and a positive composition of
$s$ to obtain this formula. Choosing one representative of each torus tie
injects the torus shell into the lattice shell. For $d=3$, $N_3(s)=4s^2+2$;
hence the distinct-site row sum is at most
\[
A\left[\frac{4q(1+q)}{(1-q)^3}+\frac{2q}{1-q}\right].
\]
In every fixed dimension~\eqref{eq:shell} grows polynomially, so the
exponential bound is uniformly summable. This is the precise correlation
meaning of the ceilings, not a numerical prediction for $T_c$.

\section{Certificate, independent checks and limitations}\label{sec:certificate}

The supplement contains a standalone Python standard-library program
\texttt{replay.py}. Its default invocation certifies both cubic witnesses and
the general-dimensional majorant. The option \texttt{--verify} additionally
runs the two complete cubic enumerators, the $N=7,8$ lift comparison, the
two-dimensional count guards, recursive edge-colouring checks and independent
bond interval arithmetic. The option \texttt{--mutants} requires all 19
specified corruptions to fail with exit code one and a stated rejecting guard.
The harness itself returns zero only when all are rejected. There is no
unconditional corrupted-input failure branch.

For $x\ge0$, take $m> x$ an integer, $z=x/m<1$, and use a positive Taylor sum
with the next-term geometric upper remainder for $e^z$, then raise both
endpoints to the integer power $m$. For $x<0$, reciprocal bounds on $e^{-x}$
reverse the endpoints. The outward dyadic implementation rounds after each
arithmetic operation. The independent implementation uses unrounded rational
sums with 24 terms. Its square-root bounds follow from integer-square-root
floor inequalities on a fixed rational grid. For $B>1$, let
$u=(B-1)/(B+1)\in(0,1)$ and use
\[
\ln B=2\sum_{j\ge0}\frac{u^{2j+1}}{2j+1},
\]
with a positive lower partial sum and a geometric upper tail.
All fractional root comparisons use~\eqref{eq:rootcomparison}.
No floating-point number decides acceptance; floating-point output is merely
a calculated diagnostic display of already certified rational bounds.

The mutations cover the cubic constants, odd colour count, bond scale and
bound, claw degree, rooted one-edge count, Raney coefficient and arity,
bridge multiplicity, source scale and Jacobian sign, decay ratio, KP budget,
pair prefactor, both general-dimensional coefficients, and two explicit
$Cy\ge1$ violations. The supplement records its command syntax and SHA-256
byte custody. The independently implemented enumerators and interval routines
are part of the same replay, not merely unexecuted alternative formulae.

The scales were selected by numerical exploration, then pinned and certified.
No global continuous-parameter optimum is claimed. The constants are rigorous
clustering ceilings, not fitted or measured temperatures, and are not expected
to locate a physical transition sharply. The $O(d^3)$ theorem removes the
$O(d^4)$ growth of the earlier hypercubic estimate, but is not an optimal
asymptotic-dimensionality assertion. The local chronological oscillator and
pair-source trace lemmas remain the shared analytical ingredients of the
released proof~\cite{rensSharp}; the present paper supplies the full
additional geometry, source and dimension-dependent estimates.

\section*{AI assistance}
The analytical results, the design of the estimates, and the final text are the
author's responsibility. Independent proof reviews of the argument and of the
certificate program were carried out by separate AI systems, whose comments led to
several corrections recorded in the present version. The manuscript, its LaTeX
source, and the certificate program were drafted and edited with AI assistance under
the author's direction and review.
\addcontentsline{toc}{section}{AI assistance}

This preprint has not been peer reviewed.

\appendix
\section{The chronological oscillator identity}\label{app:oscillator}

For a single oscillator, let $D(c\eta,u)=
\exp(c\eta e^{\Omega u}a^\dagger-c\eta e^{-\Omega u}a)$, with
$\eta^2=U/(2\Omega)$. At times $0\le u_1<\cdots<u_m<\beta$, chronological
ordering means the latest factor is on the left:
\[
\mathcal T\prod_jD(c_j\eta,u_j)
=D(c_m\eta,u_m)\cdots D(c_1\eta,u_1).
\]
The jumps are neutral, $\sum_jc_j=0$, including physical pair-source jumps.
Scalar BCH and thermal normal ordering give
\begin{equation}\label{eq:oscidentity}
\begin{split}
\log\Tr_{\mathrm{ph},0}\!\left(\mathcal T\prod_jD(c_j\eta,u_j)\right)
={}&-\frac{\eta^2}{2}\coth(\beta\Omega/2)\left(\sum_jc_j\right)^2\\
&-\eta^2\sum_{i<j}c_ic_j
\frac{\cosh(\beta\Omega/2-\Omega(u_j-u_i))
-\cosh(\beta\Omega/2)}{\sinh(\beta\Omega/2)}.
\end{split}
\end{equation}
The trace is normalized by the free-oscillator partition function.
For neutral jumps the first term vanishes. Two integrations by parts for the
periodic step occupation identify the second term with $-\Phi_\beta[n]$ in
\eqref{eq:phi}. The positivity follows from the squared occupation difference
and the positive kernel~\eqref{eq:kernel}, not from commuting chronological
displacements. This identity includes all thermal winding and every
finite positive frequency; no anti-adiabatic limit or low-frequency
approximation is taken.

\bibliographystyle{plain}
\bibliography{refs}
\end{document}
